\section{Armor}
\subsection{Introduction}
Armor is wearable cover - very convenient for anyone not looking to die.\\
Every part of the body can be covered once, meaning that armor pieces which overlap or even cover the same locations altogether cannot be combined.

\subsection{Casual wear}
This section contains clothing that is slightly armored but looks inconspicuous.\par
\begin{multicols}{2}
\luaimport{equipment/armor-casual.csvin}{armor.tpl}{equipment/casualarmor}
\end{multicols}

\subsection{Helmets}
A section on protective head wear, armoring the arguably most vulnerable body location.\par
\begin{multicols}{2}
\luaimport{equipment/armor-helmets.csvin}{armor.tpl}{equipment/helmets}
\end{multicols}

\subsection{Body Armor}
Armor pieces below are about protecting the largest location of the human body.\par
\begin{multicols}{2}
\luaimport{equipment/armor-body.csvin}{armor.tpl}{equipment/bodyarmor}
\end{multicols}

\subsection{Leg Armor}
This section is the result of a fine balance act between all-important movement and protection.\par
\begin{multicols}{2}
\luaimport{equipment/armor-leg.csvin}{armor.tpl}{equipment/legarmor}
\end{multicols}

\subsection{Mods}
\label{subsec:armormod}
In the following section there are lists for all the modifications that can be added to armor pieces. The amount of mods that can be added to any one armor piece or combination is given at the piece of armor itself.
\par
\vspace{2mm}
\begin{multicols}{2}
\luaimport{equipment/armormods.csv}{armormod.tpl}{equipment/armormods}
\end{multicols}

\section{PES}
Personal energy shields are widely available
	and protect the body without being cumbersome,
	large or particularly heavy.
They behave like a non-Newtonian fluid:
	slow, low-energy attacks (i.e. melee attacks or primitive ranged attacks) are not blocked.
\par%
\vspace*{3mm}
\begin{multicols}{2}
\luaimport{equipment/pes.csv}{pes.tpl}{pes}
\end{multicols}
